% !TeX root = ../papers/nc-p3.tex
\section{三面の一致する量子点と九係数の中心一次拡大}
\label{sec:quantum-triple-face}
四つの向き付き面値のうちちょうど三つが一致するとする。残る面の添字を0に取る。
共通値を $Q$ とすれば $(h_0,h_1,h_2,h_3)=(Q^{-3},Q,Q,Q)$ で $Q^4\ne1$ である。
向きを反転すると $Q$ と $Q^{-1}$ が交換されるので各背景について一方を固定する。
文字を $t,x_1,x_2,x_3$ とし、以下 $(a,b,c)$ は $(1,2,3)$ の巡回置換を走る。
\begin{theorem}\label{thm:quantum-triple-face}
この量子点を軌道閉包に持つ全窓の表示は
\[
 [t,x_a]=0,\qquad x_bx_c-Qx_cx_b=A_a x_a^2+B_a tx_a+D_a t^2,
 \qquad A_1A_2A_3=0
\]
で尽くされる。全て標準Koszul左右AS正則で最小周期は1である。
$C_{a,0}=A_a,C_{a,1}=B_a,C_{a,2}=D_a$ に指標
$u_{a,s}=se_0+(2-s)e_a-e_b-e_c$ を与える。
二点が添字保持同型となる必要十分条件は、三添字の巡回置換後に非零台 $S$ が同じで、
$\mathbb Z^S\to\mathbb Z^4,\ m\mapsto\sum m_{a,s}u_{a,s}$ の整数核の各基底元で
$\prod C_{a,s}^{m_{a,s}}$ が一致することである。極大な三つの八次元座標平面は巡回で置換され、
指定を忘れた一般同型類の次元は5である。ある点の退化は、非零台のうち余指標の値が
零の項を残し正の項を消す操作と、上記の同型だけで得られる。
ここでの退化は同じ背景のファイバー内で取る。
\end{theorem}
\begin{proof}
$t$ を中心とする量子表示は定理\ref{thm:quantum-indexed-all}の三角積から得る。
各行積は1で、面値1は高々一つである。前節の接計算で、非零横断重みは
共通値の三面から二面を選んだ九つだけで、各重複度は高々1である。
全ての相補対に $t$ が入るため位置モードは1である。

テンソル積記号を一部省略して $L_a=t x_a-x_a t$, $R_a=x_bx_c-Qx_cx_b$,
$F_a=A_ax_a^2+B_a(t x_a+x_a t)/2+D_at^2$ を二文字空間に取る。四テンソル
\[
 W(C)=\sum_a \bigl(L_a\otimes(R_a-F_a)+(R_a-F_a)\otimes L_a\bigr)
\]
は一位置回転で符号を変える。背景部分は順列語の係数を巡回して分かる。
補正部分は各 $s$ で四語の符号付き巡回和の定数倍となる。
$L_a,R_a-F_a$ の六本は独立である。実際、相異なる非中心文字の積への射影で
後の三本の係数が消え、前の三本は別々の文字対にある。
従って両隣接切断は全係数で階数6、関係は表示の六本である。
九つの補正テンソルは非零で相異なる指標を持ち、負の添字が二つで
軌道接に入らない。従ってこれらは全非零横断接方向の基底である。

安定化群は共通対角トーラスと純巡回置換で生成される。
奇置換なら $Q=Q^{-1}$ を要するので起こらない。この群は上のアフィン九平面を保つ。
この平面を含む線型横断面を前節と同様に選ぶ。各座標への線型射影は平面上で $C_{a,s}$ と
一致するようトーラス同変に選べる。切片の完備局所環はこれらと重み零の生成元で位相的に生成される。
零ファイバーの点が余指標で背景へ流れるなら、九座標の非零台はその正半空間に含まれる。
同じ九値の平面上の点も同じ余指標で流れるので、両弧上で重み零生成元は背景での定数値に等しく、
九座標は同じ単項式となる。従って弧は一致する。飽和エタール切片による
全軌道の同定は定理\ref{thm:quantum-attraction-fiber}の証明と同じである。

$u_{1,0}+u_{2,0}+u_{3,0}=0$ なので三つの $A_a$ が同時に非零なら正半空間に入らない。
逆に消える平方係数の添字 $a$ に余重み $-2$、他の非中心二文字に1を取り、
$t$ に十分大きな正整数を与えれば、$a$ 以外の平方と六つの $B,D$ が全て正になる。
従って零ファイバーはちょうど三平面の和である。各点の一径数退化に
命題\ref{prop:unrestricted-good-reduction}を適用すれば全次数の正則性が従う。
表示は全位置で同一である。残る同型・面による退化は対角トーラス軌道の指標格子で決まる。
八座標の台では $u_{a,2}-u_{a,1}=e_0-e_a$ が三つとも取れるため
指標像は全格子であり、台を保つ非自明な巡回もない。一般次元は $8-3=5$ となる。
\end{proof}

\begin{proposition}\label{prop:triple-face-all-regularity}
上記の九係数表示で平方積 $P=A_1A_2A_3\ne0$ としても、二つの等式
$Q^3=-1,\ P=-1$ が同時に成立する場合を除き、全て標準Koszul左右AS正則である。
二等式が成立する場合は $B,D$ によらず指数成長する。
\end{proposition}
\begin{proof}
$t$ の余重みのみを1に取ると $B,D$ が消え、平方項は全て残る。
極限代数は $S[t]$ である。$S$ は対角文字変更によって三平方係数を
共通値 $r$ ($r^3=P$) に揃えた三生成Sklyanin代数である。実際、文字倍率
$d_a$ による $A_a$ の倍率は $d_a^3/(d_1d_2d_3)$ またはその逆で、積1の全組を取る。
通常の巡回関係 $p x_bx_c+q x_cx_b+s x_a^2$ のパラメータは $(p,q,s)=(1,-Q,-r)$ となる。

三生成Sklyanin代数のKoszul性とHilbert級数の定理
（Iyudu--Shkarin, \textit{Three dimensional Sklyanin algebras and Gr\"obner bases}）を用いる。
三係数が非零で三乗が全て等しい場合を除くと、$S$ はKoszulで級数 $(1-z)^{-3}$ を持つ。
左右AS性もここで直接読める。双対の次数は $1,3,3,1$ で、三つの関係 $r_a$ に対して
$\sum_a x_a\otimes r_a=\sum_a r_a\otimes x_a$ は非零な三次交叉の生成元である。
従って一次と二次の両向きの積対合は完全である。$S[t]$ はKoszulで双対は
$S^!$ と一生成外積代数の符号付きテンソル積となる。上の余指標と
命題\ref{prop:unrestricted-good-reduction}で $B,D$ を任意に戻せる。

例外では $A/(t)=S$ の成長を直接計算できる。三係数の三乗が等しいとき、
$x_j$ を $\epsilon^j$ 倍する対角自己同型（$\epsilon^3=1$）による次数twistは
関係の三係数を共通倍率を除き $(\epsilon^{-1}p,\epsilon q,s)$ に変える。
適切な $\epsilon$ で最初の二係数を揃え、先の対角文字変更で第三も揃える。
$(1,1,1)$ の関係空間は $X_i=\sum_{j=1}^3\omega^{ij}x_j$ ($i=0,1,2$、$\omega$ は原始三乗根)
の三平方が張る。実際、$b+c\equiv2a\pmod3$ で各平方は関係のFourier和である。
同じ文字の連続しない全語が基底となり、次数 $n\ge1$ の次元は $3\cdot2^{n-1}$。
twistは次数次元を保つので元の代数にも指数成長する商がある。
\end{proof}
